\documentclass[11pt,a4paper]{article}
\usepackage[margin=23mm]{geometry}
\usepackage{amsmath,amssymb,amsthm}
\setlength{\emergencystretch}{2em}
\newtheorem{theorem}{Theorem}
\newcommand{\gen}[1]{\langle#1\rangle}
\newcommand{\Alt}{\operatorname{Alt}}
\title{A final double transposition before a singleton attachment}
\author{Human Compression project: L044 continuation}
\date{}
\begin{document}\maketitle
Let $G=\gen{\rho_0,\ldots,\rho_{r-1}}\le\Alt(\Omega)$
be transitive with a string C-marking, where $r\ge4$.
Write
\[
 s=\rho_{r-1},\quad t=\rho_{r-2},\quad u=\rho_{r-3},
 \qquad P=G_{\le r-2},\quad J=G_{\le r-3}.
\]
Assume $P=\Alt(X)$ is the actual natural alternating group
on $X=\Omega\setminus\{z\}$. Assume also that $J$ has an
orbit $C$ of size at least four and contains the pure
natural subgroup $\Alt(C)$.

\begin{theorem}
If $t$ is a double transposition, the actual last three
marks form the six-point $A_5$ marking of source Table~2(3),
read from the terminal end. More precisely, there are six
distinct points $z,a,e,e',c,d$ such that
\[
 s=(z\ a)(e\ e'),\qquad
 t=(a\ e)(e'\ c),\qquad
 u=(e\ e')(c\ d),                                  \tag{1}
\]
and these three marks fix every other point.
No numerical rank bound or smaller-degree induction is
required for this implication.
\end{theorem}
\begin{proof}
The involution $s$ centralizes $J$. Transitivity of $G$
forces $a=z^s\ne z$. Since $J$ fixes $z$ and commutes with
$s$, it fixes $a$ as well. The centralizer of the natural
$\Alt(C)$ fixes $C$ pointwise when $|C|\ge4$; hence $s$
fixes $C$, and $a\notin C$.

The two edges of $t$ must connect all the $J$-orbits on
$X$, because $P=\gen{J,t}$ is transitive there. There are
therefore at most three such orbits. If all orbits other
than $C$ were singletons, $s$ would be a nonidentity even
involution on at most three points, which is impossible.
Thus the orbits are exactly $C$, $\{a\}$ and one nontrivial
orbit $E$. The involution $s$ has the form
\[
                         s=(z\ a)\delta,
\]
where $\delta$ is odd, preserves $E$, and centralizes the
transitive action of $J$ on $E$. Such a centralizing
nonidentity involution is free: a fixed point would, by
transitivity and commutation, force every point fixed.
Write $|E|=2m$; oddness of $\delta$ says that $m$ is odd.

Contracting the three $J$-orbits, the two edges of $t$
give a tree. They cannot both meet the singleton $a$,
since $t$ is an involution. There are precisely two forms:
\[
 \begin{array}{ll}
 \text{I:}&t=(a\ c)(e\ c'),\quad c,c'\in C,\ c\ne c',\ e\in E;\\
 \text{II:}&t=(a\ e)(f\ c),\quad c\in C,\ e,f\in E,\ e\ne f.
 \end{array}                                             \tag{2}
\]
Put $e'=e^\delta$. Full IP gives
\[
          \gen{s,t}\cap P=\gen t.                        \tag{3}
\]
Every element of $\gen{s,t}$ fixing $z$ is even and
supported on $X$, so belongs to the left side of (3).

In form I, $st$ is a three-cycle on each of
$\{z,a,c\}$ and $\{c',e,e'\}$. On the remaining points
it is the product of the other $m-1$ pairs of $\delta$.
Thus $(st)^3$ is exactly that remaining product. It fixes
$z$ and $a$, so (3) forces it to be the identity; $t$
moves $a$. Hence $m=1$.

In form II with $f=e'$, the product $st$ is a five-cycle
on $\{z,a,e,e',c\}$ and is $\delta$ on the remaining
$m-1$ pairs. Now $(st)^5$ and (3) again give $m=1$.
If $f\ne e'$, the product $st$ is a four-cycle on
$\{z,a,e,e'\}$ and a three-cycle on
$\{f,f^\delta,c\}$. Its fourth power fixes $z$ but
retains that three-cycle. It cannot belong to $\gen t$,
contrary to (3). This proves, in every surviving case,
\[
                      E=\{e,e'\},\qquad\delta=(e\ e').   \tag{4}
\]

Set $J'=G_{\le r-4}$. It fixes $a$ and centralizes $t$.
Because $J'$ preserves $E$ and $t$ carries a point of
$E$ outside $E$, that point is fixed by $J'$: its image
must lie both in $E$ and in the inverse $t$-image of
the other $J$-orbit. That intersection is a singleton
in either form of (2). Hence $J'$ fixes both points of
$E$. Transitivity of $J$ on $E$ then forces $u$ to swap
$e,e'$.

Consider form II. Commutation with $t=(a\ e)(e'\ c)$
now also forces $J'$ to fix $c$. Since $J$ is transitive
on $C$, $u$ moves $c$ to a distinct point $d\in C$.
Consequently
\[
                     u=(e\ e')(c\ d)\epsilon,
\]
where $\epsilon$ is an even involution supported on
$C\setminus\{c,d\}$. On the six points in (1), the
three core permutations are literally the accepted
Table~2(3) marking, with labels $1,\ldots,6$ corresponding
to $z,a,e,e',c,d$. Their group is the perfect group $A_5$.

Let $Q=\gen{s,t,u}$. Its only exterior action is
$\gen\epsilon$. If $\epsilon\ne1$ and projection to
the six-point core were injective, its exterior action
would define a nontrivial homomorphism $A_5\to C_2$:
the first two core marks map to $1$, and the third to
$\epsilon$. Perfection excludes this. Thus
$\epsilon\in Q$. As $\epsilon\in\Alt(C)\le J$, full IP
puts it in $Q\cap J=\gen u$. It fixes $E$, whereas $u$
swaps $E$, so $\epsilon=1$. We obtain exactly (1).

It remains to exclude form I. Here commutation with
$t=(a\ c)(e\ c')$ forces $J'$ to fix both $c,c'$.
Neither can be fixed by $u$, nor can $u$ exchange them:
either event would create a proper $J$-orbit inside $C$.
Thus $u$ sends them to distinct $d,d'\in C\setminus\{c,c'\}$,
and
\[
                    u=(e\ e')(c\ d)(c'\ d')\epsilon,
\]
where $\epsilon$ is an odd involution supported on
$C\setminus\{c,c',d,d'\}$. It is disjoint from the two
four-point sets
\[
                 A=\{z,a,c,d\},\qquad B=\{e,e',c',d'\}.
\]
The core projections of $Q$ on $A$ and $B$ are each
$S_4$. On $A$ their three marked restrictions are the
successive path transpositions $(z\ a),(a\ c),(c\ d)$.
On $B$ they are $(e\ e'),(e\ c'),(e\ e')(c'\ d')$;
the first and third give $(c'\ d')$, completing a path.

The word $(tu)^4$ is a pure three-cycle on $A$:
$tu$ has orders three and four on $A$ and $B$, respectively,
and its exterior factor $\epsilon$ disappears in the fourth
power. Likewise $(stu)^4$ is a pure three-cycle on $B$:
the core orders are four and three, respectively.
Conjugating the first pure three-cycle by $Q$ supplies
all three-cycles on $A$, since its projection there is
$S_4$; conjugating the second does the same on $B$.
Therefore $Q$ contains the independent pure subgroups
$\Alt(A)$ and $\Alt(B)$.

The group $A_4$ is transitive on the six unordered pairs
of its four points. Conjugate $s=(z\ a)(e\ e')$ by an
element of $\Alt(A)\times\Alt(B)$ sending these pairs
to $\{c,d\}$ and $\{c',d'\}$. This gives
\[
                     v=(c\ d)(c'\ d')\in Q.
\]
But $v\in\Alt(C)\le J$, while $Q\cap J=\gen u$.
The nonidentity $v$ fixes $E$ and $u$ does not. This
contradiction eliminates form I and completes the proof.
\end{proof}

\paragraph{Attribution and boundary.}
The literal six-point $A_5$ marking and its identification are
borrowed from the accepted Table~2 recovery, with its complete
census and source transcription retained as required inputs.
The pure prefix alternating kernel is an explicit hypothesis;
its acquired proof and root's initial kernel suggestion remain
paid where this theorem is applied. The new argument is the
orbit-tree reduction, the two-generator stabilizer test, and
the two pure $A_4$ factors excluding the other attachment.
The last-prefix double-transposition hypothesis has not been
proved for every remaining odd-degree configuration. Even when
it holds, the conclusion is an actual opposite head, not an
exclusion of all opposite heads or a complete L044 recovery.
No new computation, formal result, original or question credit
is claimed.
\end{document}
